\documentclass[border=8pt]{standalone}
\usepackage[utf8]{inputenc}
\usepackage[T1]{fontenc}
\usepackage[spanish,mexico]{babel}
\usepackage{tikz}
\usetikzlibrary{positioning,arrows.meta,fit,backgrounds,calc,shapes.geometric}

\hyphenpenalty=10000
\exhyphenpenalty=10000

% --- figura de actor UML ---
\tikzset{
  actorpic/.pic={
    \draw[black!65] (0,0) circle (2.1mm);
    \draw[black!65] (0,-2.1mm) -- (0,-7.5mm);
    \draw[black!65] (-3mm,-4.0mm) -- (3mm,-4.0mm);
    \draw[black!65] (0,-7.5mm) -- (-2.6mm,-12mm);
    \draw[black!65] (0,-7.5mm) -- (2.6mm,-12mm);
  }
}

\begin{document}
\begin{tikzpicture}[
  font=\small,
  uc/.style={ellipse, draw=black!55, fill=white, align=center,
             minimum width=40mm, minimum height=10mm, inner sep=1.5pt, font=\scriptsize},
  lnk/.style={draw=black!50},
  ext/.style={-{Latex[length=1.8mm]}, draw=black!45, dashed}
]

% ---------- casos de uso ----------
\node[uc] (u1) {Iniciar sesión};
\node[font=\tiny\itshape, text=black!55, below=1.2mm of u1]
  (nota) {Todos los demás casos de uso lo incluyen};

\node[uc, below=5mm of nota]  (u2)  {Registrar solicitud\\y obtener folio};
\node[uc, below=8mm of u2]    (u3)  {Cargar documento\\y sus metadatos};
\node[uc, below=8mm of u3]    (u4)  {Turnar expediente\\a otra oficina};
\node[uc, below=8mm of u4]    (u5)  {Registrar acuse\\de recepción};
\node[uc, below=8mm of u5]    (u6)  {Resguardar o cancelar\\expediente};
\node[uc, below=8mm of u6]    (u7)  {Consultar expediente};
\node[uc, below=8mm of u7]    (u8)  {Visualizar documento\\en el visor};
\node[uc, below=8mm of u8]    (u9)  {Consultar bitácora};
\node[uc, below=8mm of u9]    (u10) {Generar reportes\\e indicadores};
\node[uc, below=8mm of u10]   (u11) {Administrar usuarios,\\roles y catálogos};

% ---------- frontera del sistema ----------
\begin{scope}[on background layer]
  \node[draw=black!45, rounded corners=2pt, inner xsep=7mm, inner ysep=7mm,
        fit=(u1)(u11),
        label={[anchor=north, font=\scriptsize\itshape, text=black!55, yshift=-1.8mm]north:
        Sistema de trazabilidad documental}] (sis) {};
\end{scope}

% ---------- actores a la izquierda ----------
\coordinate (ca1) at ([xshift=-40mm]$(u2)!0.5!(u3)$);
\coordinate (ca2) at ([xshift=-40mm]$(u5)$);
\coordinate (ca3) at ([xshift=-40mm]$(u7)!0.5!(u8)$);
\foreach \c/\n/\t in {ca1/A1/{Capturista}, ca2/A2/{Responsable\\de oficina}, ca3/A3/{Consultante}}{
  \pic at (\c) {actorpic};
  \node[font=\scriptsize, align=center, text width=25mm] at ([yshift=-17.5mm]\c) {\t};
}

% ---------- actores a la derecha ----------
\coordinate (cb1) at ([xshift=40mm]$(u9)$);
\coordinate (cb2) at ([xshift=40mm]$(u10)$);
\coordinate (cb3) at ([xshift=40mm]$(u11)$);
\foreach \c/\t in {cb1/{Auditor}, cb2/{Directivo}, cb3/{Administrador\\del sistema}}{
  \pic at (\c) {actorpic};
  \node[font=\scriptsize, align=center, text width=25mm] at ([yshift=-17.5mm]\c) {\t};
}

% ---------- asociaciones ----------
\foreach \c/\u in {ca1/u2, ca1/u3, ca2/u4, ca2/u5, ca2/u6, ca3/u7, ca3/u8}
  \draw[lnk] ([yshift=-5mm]\c) -- (\u.west);
\foreach \c/\u in {cb1/u9, cb2/u10, cb3/u11}
  \draw[lnk] ([yshift=-5mm]\c) -- (\u.east);

% ---------- extension del visor ----------
\draw[ext] (u8.east) to[out=25,in=-25]
  node[midway, right=1.5pt, font=\tiny\itshape, text=black!55]{$\ll$extend$\gg$} (u7.east);

\end{tikzpicture}
\end{document}
